\clearpage
\fancyhead[L]{\small HALLASCHKA. ELEMENTA ECLIPSIUM.}
\SourceLeaf{624}
\begin{center}\small Göttingische gelehrte Anzeigen. Stück 17. Seite 167..168. 1820 Januar 29.\\[3pt]\rule{.72\linewidth}{.3pt}\end{center}

\emph{Elementa eclipsium quas patitur tellus, luna eam inter et solem versante, ab a. {\upshape 1816} usque ad a. {\upshape 1860}, ex tabulis astronomicis recentissime conditis et calculo parallactico deducta, typo ecliptico et tabulis projectionis geographicis collustraia, a} \textsc{Cassiano Hallaschka} Prag 1816. Typis Theophili Haase. 107 Seiten in 4. nebst 20 Kupfertafeln.

Bei weitem den grössten Theil dieser Schrift nimmt eine an den Sonnenfinsternissen von 1816 und 1818 practisch gelehrte und bei ersterer mit dem ganz ausführlichen Rechnungsdetail begleitete Anleitung zur parallactischen Berechnung und orthographischen Entwerfung der Sonnenfinsternisse ein. Für die ersten Anfänger kann diese, nur auf den gewöhnlichen kunstlosen Methoden beruhende, Anweisung ihren Nutzen haben, in welcher Rücksicht jedoch der Verf. bei den vorausgeschickten Vorschriften etwas mehr Sorgfalt auf einen deutlichen und richtigen Vortrag hätte verwenden sollen. Auf den letzten 18 Seiten gibt der Verf. noch die Hauptmomente der in Brünn vom Jahr 1818 bis 1860 sichtbaren Finsternisse an, so wie die beigefügten Kupfertafeln die allgemeine Darstellung derselben für die ganze Erde liefern. Es sind darunter drei partiale, fünf ringförmige und vier totale. Für den Horizont von Brünn ist die vom 8ten Juli 1842 die grösste, und die Grösse der Verfinsterung beträgt daselbst 11 Zoll $13'$. Auf die Zeichnung selbst scheint nicht durchaus die grösste Sorgfalt gewandt zu sein: so finden wir bei der erwähnten Finsterniss die Punkte, wo Anfang oder Ende bei Aufgang oder Untergang der Sonne eintritt, als Eine zusammenhängende Linie gezeichnet, da sie doch nothwendig zwei getrennte und durch die Linie des Contactus limborum solis australis et lunae borealis verbundene Linien bilden mussten. Die Rechnungen sind nach \textsc{Triesnecker}'s Mondstafeln in den Wiener Ephemeriden für 1803 und nach handschriftlichen Sonnentafeln desselben Astronomen geführt. Der erste Meridian ist nicht, wie es sonst üblich ist, $20^\circ$ sondern $19^\circ\,53'\,45''$ westlich von Paris gesetzt.
